\section{Generative UI}

\label{sec:genui}

\subsection{Component System}

Hanzo Search renders answers using a library of 12 generative UI components, selected based on query intent:

\begin{table}[H]
\centering
\small
\begin{tabular}{lll}
\toprule
\textbf{Component} & \textbf{Intent} & \textbf{Description} \\
\midrule
AnswerCard & Factual & Concise answer with sources \\
StepList & How-to & Numbered steps with details \\
CompareTable & Comparison & Side-by-side feature table \\
Timeline & Historical & Chronological event display \\
CodeBlock & Technical & Syntax-highlighted code \\
Calculator & Computational & Interactive calculator \\
Chart & Data & Bar/line/pie chart rendering \\
MapView & Location & Geographic visualization \\
ProductCard & Shopping & Product info with pricing \\
PerspectiveCards & Opinion & Multiple viewpoint display \\
LinkCard & Navigational & Direct link with preview \\
FreeForm & Creative & Unrestricted rich text \\
\bottomrule
\end{tabular}
\caption{Generative UI component library.}
\label{tab:components}
\end{table}

\subsection{Component Selection}

The intent classifier (\S\ref{sec:query}) determines the primary component. A secondary pass by the LLM may add supplementary components (e.g., a factual answer with a chart).

\begin{algorithm}[H]
\caption{Generative UI Rendering}
\label{alg:genui}
\begin{algorithmic}[1]
\Require Answer $a$, citations $\mathcal{C}$, intent $c$, data $\mathcal{D}$
\State $\text{primary} \gets \text{IntentToComponent}(c)$
\State $\text{structured} \gets \text{LLM.Structure}(a, \text{primary})$
\State \Comment{Extract structured data for rich components}
\If{$c = \text{comparison}$}
    \State $\text{table} \gets \text{LLM.ExtractTable}(a)$
    \State $\text{structured.table} \gets \text{table}$
\ElsIf{$c = \text{computational}$}
    \State $\text{chart\_data} \gets \text{LLM.ExtractData}(a, \mathcal{D})$
    \State $\text{structured.chart} \gets \text{chart\_data}$
\EndIf
\State $\text{UI} \gets \text{React.Render}(\text{structured}, \mathcal{C})$
\State \Return $\text{UI}$
\end{algorithmic}
\end{algorithm}

\subsection{Streaming Architecture}

Hanzo Search streams both text and UI components to the user using Server-Sent Events (SSE). The streaming protocol supports three event types:

\begin{enumerate}[leftmargin=1.4em]
    \item \textbf{text\_delta}: Incremental text tokens for the answer body.
    \item \textbf{component}: A complete UI component (table, chart) sent when ready.
    \item \textbf{citation}: A source citation with URL, title, and snippet.
\end{enumerate}

This enables progressive rendering: users see the text answer streaming in real-time while structured components (tables, charts) appear as soon as the relevant data is extracted.

\subsection{User Engagement Results}

A/B testing with 50K users over 30 days (25K per condition):

\begin{table}[H]
\centering
\small
\begin{tabular}{lcc}
\toprule
\textbf{Metric} & \textbf{Text-only} & \textbf{Generative UI} \\
\midrule
Time on page & 42s & 67s (+59\%) \\
Follow-up queries & 2.3 & 1.2 ($-$48\%) \\
Task completion & 61\% & 82\% (+34\%) \\
User satisfaction & 3.8/5 & 4.4/5 (+16\%) \\
Click-through to sources & 34\% & 52\% (+53\%) \\
\bottomrule
\end{tabular}
\caption{A/B test results: generative UI vs text-only answers.}
\label{tab:ab_test}
\end{table}

The reduction in follow-up queries ($-$48\%) indicates that generative UI answers are more self-contained, reducing the need for refinement.

